\section{Rangkuman}
\label{sec:rangkuman-bab6}

Pada bab ini telah dibahas:

\begin{itemize}
    \item Register 16 bit hanya menjangkau 64 KB, sedangkan bus
          alamat 8086 selebar 20 bit menjangkau 1 MB. Segmentasi
          menjembatani kedua sisi tersebut.

    \item Alamat logis ditulis sebagai pasangan
          \eng{segment:offset}, misalnya \code{1234h:5678h}.
          Alamat fisik adalah alamat 20 bit yang dikirim ke bus
          alamat.

    \item Alamat fisik dihitung dengan rumus
          $(\text{segmen} \times \hex{10}) + \text{offset}$.
          Pengali \hex{10} sama dengan 16 desimal.

    \item Memori dibagi menjadi paragraf berukuran 16 byte. Nilai
          segmen menyatakan nomor paragraf, sehingga alamat dasar
          segmen selalu kelipatan 16.

    \item Satu segmen menjangkau 64 KB, dari offset \hex{0000}
          sampai \hex{FFFF}. Enam belas segmen menutupi 1 MB.

    \item \eng{Overlapping} segmen membuat satu alamat fisik dapat
          dinyatakan oleh banyak pasangan \eng{segment:offset}.
          Jumlah pasangan untuk alamat $F$ di bawah \hex{10000}
          adalah $\lfloor F/16 \rfloor + 1$, dan maksimumnya 4.096.

    \item Memori 1 MB terbagi menjadi memori konvensional 640 KB,
          buffer video, ROM adaptor, dan ROM BIOS.

    \item Alamat fisik dapat diverifikasi dengan DEBUG atau
          emu8086. Perintah \lstasm|D segmen:offset| menampilkan
          alamat logis; alamat fisiknya dihitung dengan rumus.
\end{itemize}

\begin{checklistbox}
\vspace{1.1em}
Tandai setiap butir yang sudah Anda kuasai.

\begin{itemize}
    \item[\checkbox] Saya dapat menjelaskan alasan segmentasi pada
          8086.
    \item[\checkbox] Saya dapat membedakan alamat logis dan alamat
          fisik.
    \item[\checkbox] Saya dapat menghitung alamat fisik dari
          \eng{segment:offset}.
    \item[\checkbox] Saya dapat menjelaskan paragraf 16 byte dan
          segmen 64 KB.
    \item[\checkbox] Saya dapat menentukan banyaknya pasangan yang
          menunjuk satu alamat fisik.
    \item[\checkbox] Saya dapat menyebutkan pembagian memori 1 MB.
    \item[\checkbox] Saya dapat memverifikasi alamat dengan DEBUG
          atau emu8086.
\end{itemize}

\vspace{0.4em}
Bila masih ada butir yang belum bertanda, ulangi bagian yang
bersesuaian sebelum melanjutkan ke bab berikutnya.
\end{checklistbox}

\begin{infobox}
\vspace{1.1em}
\textbf{Rujukan bab.} Materi pada bab ini diringkas dari
\cite{gunawan2019}, Bagian 2.1.1, hlm.~13--15, dan \cite{sto2001},
BAB II, hlm.~9--13.
\end{infobox}
